\subsection{连续实值函数的扩张}

设 X 与 Y 是两个拓扑空间，\(A \neq X\) 是 X 的闭子集。若 f 是 A 到 Y 中的连续映射，未必总能把 f 扩张为整个 X 到 Y 中的连续映射。当 \(Y = \overline{R}\) 时，这种扩张的可能性由下面的定理确定：

定理 2（乌雷松）。公理 \((\mathbf{O}_{\mathbf{V}})\) 等价于下列性质：\((\mathbf{O}_{\mathbf{V}}^{\prime \prime \prime})\) 给定 X 的任一闭子集 A 与定义在 A 上的任一连续实值函数 \(f\) （有限的或非有限的），存在一个扩张 \(g\) 把 \(f\) 延拓到整个空间 X，且是 X 到 \(\overline{\mathbf{R}}\) 中的连续映射。

容易看出 \((\mathbf{O}_{\mathbf{V}}^{\prime\prime})\) 蕴含 \((\mathbf{O}_{\mathbf{V}})\) ；因为若 B 与 C 是 X 的两个不相交闭子集，则在 B 上等于 o、等于 \(\mathbf{I}\) 于 \(\mathbf{C}\) 上的函数在闭集 \(\mathbf{B} \cup \mathbf{C}\) 上有定义且连续，故由 \((\mathrm{O}_{\mathbf{V}}^{\prime \prime})\) 它有连续扩张 \(f\) 到 \(\mathbf{X}\) 上。若 \(g = \inf (f^{+}, \mathbf{I})\) ，则 \(g\) 在 \(\mathbf{X}\) 上连续，取值于 \([0, \mathbf{I}]\) 且等于 \(\mathbf{0}\) 于 \(\mathbf{B}\) 上、等于 \(\mathbf{I}\) 于 \(\mathbf{C}\) 上。

我们来证明反之 \((\mathbf{O}_{\mathbf{v}})\) 蕴含 \((\mathbf{O}_{\mathbf{v}}^{\prime\prime})\) 。由于 \(\overline{R}\) 与区间 \([-1, +1]\) 同胚，只需考虑连续映射 \(f: A \to \overline{R}\) 取值于 \([-1, +1]\) 的情形。我们将通过构造 X 上的连续函数序列 \((g_{n})\) 来作出 f 到 X 的扩张 g，使得序列 \((g_{n}(x))\) 对一切 \(x \in X\) 收敛于区间 \([-1, +1]\) 的一点；这个极限按定义就是 g 在 x 处的值，并且由 \(g_{n}\) 的选法将得出函数 g 满足所要求的条件。

\(g_{n}\) 的定义基于下面的引理：

引理 1. 设 u 是 A 到 \([-1, +1]\) 中的连续映射；则存在 X 到 \([-1/3, +1/3]\) 中的连续映射 v，使得 \(|u(x) - v(x)| \leqslant 2/3\) 对一切 \(x \in A\) 成立。

设 H 是全体 \(x \in A\) 中满足 \(-\mathtt{I} \leqslant u(x) \leqslant -\mathtt{I}/3\) 者所成的集，K 是全体 \(x \in A\) 中满足 \(\mathtt{I}/3 \leqslant u(x) \leqslant \mathtt{I}\) 者所成的集；H 与 K 在 A 中从而在 X 中是闭的，且不相交；故由 \((\mathrm{O}_{\mathrm{V}})\) 存在 X 到 \([-\mathtt{I}/3, +\mathtt{I}/3]\) 中的连续映射 v，它在 H 上等于 \(-\mathtt{I}/3\) 而在 K 上等于 \(+\mathtt{I}/3\) 。映射 v 满足引理的条件。

我们现在归纳地定义函数 \(g_n\) 。对 \(u = f\) 应用引理，定义 \(g_0\) 为 \(\mathbf{X}\) 到 \([-1/3, +1/3]\) 中的连续映射，使得 \(|f(x) - g_0(x)| \leqslant 2/3\) 对一切 \(x \in \mathbf{A}\) 成立。设连续映射 \(g_n\) （\(\mathbf{X}\) 到区间

\[
[ - \mathrm{I} + (\frac {2}{3}) ^ {n + 1}, \mathrm{I} - (\frac {2}{3}) ^ {n + 1} ]
\]

已定义，使得 \(|f(x) - g_n(x)| \leqslant (\frac{2}{3})^{n+1}\) 对一切 \(x \in \mathbf{A}\) 成立。对函数 \(u(x) = (\frac{2}{3})^{n+1} (f(x) - g_n(x))\) 应用引理，可见存在连续映射 \(h_{n+1}\) （\(\mathbf{X}\) 到区间 \(\left[-\frac{2^{n+1}}{3^{n+2}}, \frac{2^{n+1}}{3^{n+2}}\right]\) 中的），使得

\[
\left| f (x) - g _ {n} (x) - h _ {n + 1} (x) \right| \leqslant (2 / 3) ^ {n + 2}
\]

对一切 \(x \in \mathbf{A}\) 成立；取 \(g_{n+1} = g_n + h_{n+1}\) 便完成了归纳，因为这个函数满足不等式 \(|g_{n+1}(x)| \leqslant 1 - (2/3)^{n+2}\) （对一切 \(x \in \mathbf{X}\) ，由 \(h_{n+1}\) 的定义）。

由这个定义推出，对 \(m \geqslant p\) 与 \(n \geqslant p\) ，我们有

\[
\left| g _ {m} (x) - g _ {n} (x) \right| \leqslant \frac {2 ^ {p + 1}}{3 ^ {p + 2}} \sum_ {k = 0} ^ {\infty} (2 / 3) ^ {k} = (2 / 3) ^ {p + 1}
\]

在每一点 \(x \in \mathbf{X}\) 处成立；因此序列 \((g_n(x))\) 对每个 \(x \in \mathbf{X}\) 是柯西序列，从而收敛于一点 \(g(x)\) （区间 \([-1, +1]\) 中的）；并且由于 \(f(x) - g_n(x)\) 对一切 \(x \in \mathbf{A}\) 而言当 \(n \to \infty\) 时趋于 o，\(g\) 是 \(f\) 到 \(\mathbf{X}\) 的扩张。因此只需再证明 \(g\) 在 \(\mathbf{X}\) 上连续。

现在设 \(x\) 是 \(\mathbf{X}\) 的任一点；于是，给定任意 \(\varepsilon > 0\) ，存在整数 \(n_0\) 使得 \(|g_m(y) - g_n(y)| \leqslant \varepsilon\) 对一切 \(y \in \mathbf{X}\) 、一切 \(m \geqslant n_0\) 与一切 \(n \geqslant n_0\) 成立；因此，令 \(m\) 趋于 \(+\infty\) ，就有

\[
\left| g (y) - g _ {n} (y) \right| \leqslant \varepsilon .
\]

设 \(\mathbf{V}\) 是 \(x\) 的一个邻域，使得 \(|g_n(y) - g_n(x)| \leqslant \varepsilon\) 对一切 \(y \in \mathbf{V}\) 成立；于是对每个 \(y \in \mathbf{V}\) 我们将有

\[
\left| g (y) - g (x) \right| \leqslant \left| g (y) - g _ {n} (y) \right| + \left| g _ {n} (y) - g _ {n} (x) \right| + \left| g (x) - g _ {n} (x) \right| \leqslant 3 \varepsilon ,
\]

这表明 g 在 x 处连续，从而完成了定理 2 的证明。（证明的最后部分以一种特殊情形用到了一致收敛的概念，我们将在第 X 章第 1 号中作一般研究。）

推论. 若 \(f\) 是定义在 \(\mathbf{A}\) 上的有限连续实值函数，则存在有限连续实值函数 \(g\) （定义在 \(\mathbf{X}\) 上的），它是 \(f\) 的扩张。

先考虑 \(f(x) \geqslant 0\) 对一切 \(x \in A\) 成立的情形；这时存在连续扩张 \(g_1\) 把 \(f\) 扩张到 \(X\) 上，取值于 \([0, +\infty]\) 。若令 \(B = \overline{g_1} (+\infty)\) ，则 \(B\) 是闭集且由假设不与 \(A\) 相交；于是函数 \(h\) ——等于 \(f\) 于 \(A\) 上、等于 \(o\) 于 \(B\) 上——是闭集 \(A \cup B\) 上的连续函数。设 \(g_2\) 是 \(h\) 到 \(X\) 上的连续扩张，它仍取值于 \([0, +\infty]\) ；则函数 \(g = \inf(g_1, g_2)\) 是 \(f\) 到 \(X\) 上的连续扩张，其值 \(\geqslant 0\) 且在 \(X\) 的每一点处有限。

要过渡到一般情形，只需注意：若 \(f\) 在 \(\mathbf{A}\) 上有限且连续，则 \(f^{+}\) 与 \(f^{-}\) 也如此；把 \(f^{+}\) 与 \(f^{-}\) 扩张到 \(\mathbf{X}\) 上，分别以有限连续函数 \(g_{1}\) 与 \(g_{2}\) 为扩张，则函数 \(g_{1} - g_{2}\) 在 \(\mathbf{X}\) 上有限且连续并扩张 \(f\) 。

注。若 X 是正规空间且 A 是 X 的闭子集，则 A 到立方体 \(K^{I}\) （§ 1，第 5 号）中的每个连续映射 f 也存在到 X 上的连续扩张；因为这时 \(f = (f_{i})_{i \in I}\) ，其中 \(f_{i}\) 是 A 到 R 中紧区间 K 中的连续映射；由于存在连续映射 \(g_{i}: X \to K\) 扩张 \(f_{i}\) ，映射 \(g = (g_{i})\) 就是 f 到 X 上的连续扩张。

